\begin{description}
\item[(T7.1)] The magnification will be given by,
\begin{align*}
m_0 &= \frac{f_{\mathrm{o}}}{f_{\mathrm{e}}} \tag*{(1.0)}\\
&= \frac{100}{1} = 100 \tag*{(1.0)}
\end{align*}
The magnification is $m_0 = 100$.

Length of the telescope will be
\begin{align*}
L_0 &= f_{\mathrm{o}} + f_{\mathrm{e}} \tag*{(1.0)}\\
&= 100 + 1 = 101\ \mathrm{cm} \tag*{(1.0)}
\end{align*}
The telescope length will be $L_0 = 101$\,cm.

\item[(T7.2)] We use the following sign convention. Lens is the origin.
Direction along the direction of light is taken as positive. The lens formula
is $\dfrac{1}{f} = \dfrac{1}{v} - \dfrac{1}{u}$ ($f$ is positive for convex
lens and negative for concave lens). Magnification is $m = \dfrac{v}{u}$.
Solutions using other sign conventions are acceptable.

Let $v$ be image distance from the Barlow lens.
\[
\frac{1}{f_{\mathrm{B}}} = \frac{1}{v} - \frac{1}{u} \tag*{(1.0)}
\]
Distance of Barlow lens, $d_{\mathrm{B}}$, before the prime focus is same as
the object distance, $u$, in this case. \hfill(1.0)
\[
\frac{1}{f_{\mathrm{B}}} = \frac{1}{v} - \frac{1}{d_B}
\]
Also, $m_{\mathrm{B}} = 2 = \dfrac{v}{u} = \dfrac{v}{d_{\mathrm{B}}}$
\hfill(0.5)
\begin{align*}
\therefore\ \frac{1}{d_{\mathrm{B}}} &= \frac{2}{v} \tag*{(1.0)}\\
\therefore\ \frac{1}{-1} &= \frac{1}{v} - \frac{2}{v}\\
-1 &= \frac{-1}{v}\\
v &= 1\ \mathrm{cm}\\
d_{\mathrm{B}} &= \frac{v}{2} = \frac{1\ \mathrm{cm}}{2} = 0.5\ \mathrm{cm}
  \tag*{(2.5)}
\end{align*}
The positive sign for $d_{\mathrm{B}}$ indicates that the Barlow lens was
introduced $\boxed{0.5\ \mathrm{cm}}$ before the prime focus.

\item[(T7.3)] The increase in the length will be,
\begin{align*}
\Delta L &= v - d_{\mathrm{B}} \tag*{(2.0)}\\
&= 1.0 - 0.5 = 0.5\ \mathrm{cm} \tag*{(2.0)}
\end{align*}
Thus, the length will be increased by $\boxed{\Delta L = 0.5\ \mathrm{cm}}$.

\item[(T7.4)] Plate scale at prime focus is given by,
\[
s = \frac{1}{f_o} = \frac{1\ \mathrm{rad}}{1\ \mathrm{m}}
 = 0.206\,265\ \mathrm{arcsec}/\mu\mathrm{m} \tag*{(2.0)}
\]
Since each pixel is $10\ \mu$m in size,
\[
s_{\mathrm{p}} = 10 \times 0.206\ \mathrm{arcsec}/\mu\mathrm{m}
 = 2.06\ \mathrm{arcsec/pixel} \tag*{(2.0)}
\]
Two stars will be separated by,
\[
n_{\mathrm{p}} = \frac{20''}{2.0600}\ \mathrm{pixels}
 \simeq 10\ \mathrm{pixels} \tag*{(2.0)}
\]
\end{description}
